\documentclass[12pt]{article}
\usepackage[T1]{fontenc}
\usepackage[utf8]{inputenc}
\usepackage{lmodern}
\usepackage{textcomp}
\usepackage{enumitem}
\usepackage{geometry}
\usepackage{graphicx}
\usepackage{amsmath, amssymb}
\geometry{margin=1in}
\begin{document}
\begin{center}
Political Science - Graduate Level Assessment 25 \
Total Marks: 40 \
Answer all questions.
\end{center}

\section*{Multiple Choice (10 marks)}
\begin{enumerate}[label=\arabic*.]
\item The dominant course for foreign policy throughout most of American history can be categorized as
\begin{enumerate}[label=(\alph*)]
    \item containment.
    \item neoconservatism.
    \item isolationism.
    \item protectionism.
\end{enumerate}
\item What did Paul Kennedy argue in his book The Rise and Fall of the Great Powers?
\begin{enumerate}[label=(\alph*)]
    \item All of the world's leading economies were declining due to low growth and inflation
    \item The United States could no longer remain a superpower and was in decline
    \item The soft power of the United States would allow it to avoid decline
    \item The rise of Japan had been exaggerated
\end{enumerate}
\item Which of the following ideas have become mainstreamed within human security since the Human Development Report?
\begin{enumerate}[label=(\alph*)]
    \item Limiting security analysis to military engagement is too restrictive.
    \item Security is a contestable concept.
    \item Focusing security only on violent conflicts is too restrictive.
    \item All of these options.
\end{enumerate}
\item In American government, the power to declare war rests with
\begin{enumerate}[label=(\alph*)]
    \item the president of the United States.
    \item the secretary of defense.
    \item the chairman of the Joint Chiefs of Staff.
    \item Congress.
\end{enumerate}
\item Within Critical Security Studies, what is the referent object of security?
\begin{enumerate}[label=(\alph*)]
    \item There is no specific referent point.
    \item The state and state apparatus.
    \item The environment.
    \item The people and their collectives.
\end{enumerate}
\end{enumerate}

\section*{True / False (10 marks)}
\textit{Answer True or False}
\begin{enumerate}[label=\arabic*.]
\setcounter{enumi}{5}
\item True or False: Digitalised sensitive information is considered the main referent object in contemporary cyber-security frameworks.
\item True or False: Approximately 30-40 million people are estimated to be living with HIV/AIDS worldwide.
\item True or False: The 'New Populism' represents a strong internationalist sentiment advocating for increased global cooperation and alliances.
\item True or False: Idealism argues that states are the sole actors in international relations, focusing primarily on moral principles and cooperation.
\item True or False: Most states openly disclose their nuclear capabilities to promote transparency and international trust.
\end{enumerate}

\section*{Long Form Response (20 marks)}
\textit{Answer each question with a clear and complete explanation.}
\begin{enumerate}[label=\arabic*.]
\setcounter{enumi}{10}
\item Discuss how inefficient balancing or buckpassing by states can contribute to a more competitive international system, providing examples and implications for global politics.
\item Discuss the various factors that contribute to the differing relationships between U.S. and European security studies, and analyze their implications for international relations.
\end{enumerate}

\vfill
\noindent\textit{End of Paper}
\end{document}